\documentclass[12pt,a4paper,openright,twoside]{book}
\usepackage[utf8]{inputenc}
\usepackage{disi-thesis}
\usepackage{code-lstlistings}
\usepackage{notes}
\usepackage{shortcuts}
\usepackage{acro}
\usepackage{hyperref}

\hypersetup{
    colorlinks=true,
    linkcolor=cyan,
    urlcolor=blue,
    citecolor=red
}
% \hypersetup{hidelinks} % PER STAMPA

%----------------------------------------------------------------------------------------
% FRONTESPIZIO
%----------------------------------------------------------------------------------------
\school{\unibo}
\programme{Corso di Laurea in Ingegneria e Scienze Informatiche}
\title{Sicurezza Kernel-Level in ambienti Cloud-Native: il caso di studio Falco e Tetragon}
\author{Matteo Grandini}
\date{\today}
\subject{Virtualizzazione e Integrazione di Sistemi}
\supervisor{Prof. Vittorio Ghini}
% \cosupervisor{Dott. CoSupervisor 1}
% \morecosupervisor{Dott. CoSupervisor 2}
\session{II}
\academicyear{2025-2026}

%----------------------------------------------------------------------------------------
% ACRONIMI
%----------------------------------------------------------------------------------------
\acsetup{case-sensitive=false}
\DeclareAcronym{ebpf}{
    short={eBPF},
    long={extended Berkeley Packet Filter},
    first-style=short-long
}
\DeclareAcronym{bpf}{
    short={BPF},
    long={Berkeley Packet Filter}, 
    first-style=short-long
}
\DeclareAcronym{vm}{
    short={VM},
    long={Virtual Machine}
}

%----------------------------------------------------------------------------------------

\mainlinespacing{1.241}

\begin{document}

\frontmatter\frontispiece

%----------------------------------------------------------------------------------------

\begin{abstract}	
Massimo 2000 caratteri. % TODO: Scrivere verso la fine.
\end{abstract}

%----------------------------------------------------------------------------------------

{\hypersetup{hidelinks}
\tableofcontents   
% \listoffigures
}

%----------------------------------------------------------------------------------------

\mainmatter

%----------------------------------------------------------------------------------------

\chapter{Introduzione}
\label{chap:introduction}
% Context & Problem Statement: La sicurezza della runtime nei container e i limiti dell'approccio tradizionale basato solo su analisi statiche o agenti in user-space.
% La Svolta eBPF: Introduzione al tracciamento a livello kernel come nuovo paradigma di sicurezza Cloud-Native.
% Obiettivi della Tesi: Definizione dell'obiettivo scientifico: confronto tra modelli reattivi (Falco) e preventivi (Tetragon), sviluppo di una suite automatizzata di benchmarking, e valutazione di integrabilità/efficacia in diversi ambienti.
% \paragraph{Struttura della tesi}

%----------------------------------------------------------------------------------------

\chapter{Background}
\label{chap:background}
% Sicurezza Cloud-Native: Limitazioni degli agenti user-space, complessità di gestione e tuning delle regole, necessità di visibilità fine-grained.
% Fondamenti di eBPF: Meccanismi di aggancio nel kernel Linux (kprobes, uprobes, tracepoints), eBPF maps e ruolo del Verifier.

Questo capitolo fornisce il contesto teorico necessario per comprendere le motivazioni alla base della scelta di strumenti come Falco e Tetragon per la sicurezza in ambienti cloud-native. Verranno discussi i concetti chiave relativi alla sicurezza dei container e al ruolo di eBPF come tecnologia abilitante per l'osservabilità e il controllo a livello kernel.

\section{Sicurezza Cloud-Native}
\label{sec:cloud-native-security}
% TODO: aggiungere breve definizione di cloud-native
% TODO: controllare riferimenti
La sicurezza in ambienti cloud-native presenta sfide specifiche dovute all'adozione diffusa di container, orchestratori e pipeline di deployment che introducono dinamismo, multitenancy e complessità nella superficie d'attacco. Rispetto ai modelli tradizionali basati su macchine virtuali monolitiche, i workload cloud-native tendono a essere più effimeri, frammentati in più servizi e soggetti a frequenti aggiornamenti: queste caratteristiche rendono inefficace una difesa statica basata esclusivamente su immagini "trusted" o su scansioni periodiche. % da rivedere

Un problema centrale è la perdita di visibilità tra il livello applicativo e il livello di sistema: i container condividono lo stesso kernel dell'host e molte attività malevole (ad esempio, escalation di privilegio, manipolazione dei filesystem montati o esfiltrazione di dati) si manifestano attraverso syscall e comportamenti runtime che non sono immediatamente evidenti osservando solo i file binari o i log applicativi \cite{kubernetes-docs,docker-docs}.

Per questo motivo, le strategie di difesa moderne nel paradigma cloud-native si articolano su tre pilastri complementari:
\begin{itemize}
    \item \textbf{Riduzione della superficie d'attacco:} uso del principio del privilegio minimo, immutabilità delle immagini, meccanismi di namespace e cgroups per isolare risorse.
    \item \textbf{Pre-deployment security:} scansione delle immagini, policy di admission nel cluster e verifica di conformità durante la pipeline CI/CD.
    \item \textbf{Runtime security e osservabilità:} monitoraggio continuo del comportamento dei processi e delle risorse, correlazione di eventi e capacità di risposta automatica o manuale.
\end{itemize}
In particolare, l'osservabilità a livello runtime richiede strumenti in grado di collegare eventi di alto livello (ad es. richieste di rete, creazione di processi, accesso a file sensibili) con le entità del cluster (pod, container, namespace, service account). Le soluzioni user-space tradizionali possono fornire ricche informazioni contestuali, ma soffrono del fatto che molte operazioni critiche avvengono nel kernel: il passaggio ripetuto di dati tra kernel e user-space può introdurre overhead e latenza, limitando la granularità e la tempestività della rilevazione \cite{ebpf-docs}.

L'avvento di tecniche basate su eBPF ha cambiato lo scenario: eBPF consente di eseguire piccoli programmi sandboxati direttamente nel kernel per osservare eventi, filtrare e aggregare dati, e comunicare informazioni a user-space in modo efficiente \cite{linux-bpf-docs,linux-btf-docs}. Questo approccio riduce le copie dei dati e la latenza, permettendo sia una migliore visibilità sia la possibilità, in alcuni sistemi, di intervenire in-linea per bloccare o mitigare azioni pericolose.

Dal punto di vista operativo, una strategia completa di sicurezza cloud-native dovrebbe combinare controlli preventivi (policy di image scanning, admission control e least-privilege) con strumenti di detection e response runtime capaci di correlare segnali multipli e, quando necessario, applicare enforcement a bassa latenza. Strumenti come Falco e Tetragon rappresentano due approcci complementari a questa sfida: Falco si focalizza sul rilevamento e l'integrazione con workflow di alerting e risposta, mentre Tetragon sfrutta la potenza di eBPF per fornire osservabilità profonda e potenziali azioni di enforcement a livello kernel \cite{falco-docs,tetragon-home}.

% Infine, va sottolineato che l'efficacia delle soluzioni runtime dipende fortemente dal contesto operativo: regole troppo permissive generano rumore e falsi positivi, regole troppo restrittive possono interrompere workload legittimi. Pertanto la sicurezza cloud-native è tanto una questione tecnologica quanto organizzativa: richiede processi di tuning, pipeline di test per le policy e integrazione con le pratiche DevOps per ridurre il rischio operativo e mantenere la resilienza del sistema.

\section{Fondamenti di eBPF}
\label{sec:eBPF-fundamentals}
\ac{ebpf} è un'estensione del meccanismo \ac{bpf} che consente l'esecuzione di programmi sandboxed all'interno del kernel Linux. Questi programmi possono essere utilizzati per monitorare e, in alcuni casi, influenzare il comportamento del sistema operativo in tempo reale, senza modificare il codice del kernel né caricare moduli proprietari \cite{ebpf-docs,linux-bpf-docs}.

Dal punto di vista architetturale, un programma \ac{ebpf} viene caricato in risposta a specifici hook del kernel, come kprobe, uprobe, tracepoint, socket filter e programmi di rete. Quando viene attivato, il programma può leggere o scrivere dati in apposite strutture chiamate mappe, scambiare informazioni con user space e, in alcuni casi, modificare il flusso di esecuzione di un evento. La sicurezza è garantita dal verifier, un componente del kernel che analizza il programma prima dell'esecuzione per verificarne la terminazione, l'assenza di loop non controllati e l'accesso corretto alle risorse \cite{ebpf-docs,linux-bpf-docs}.

Questa architettura rende \ac{ebpf} particolarmente adatto a scenari in cui è necessario ottenere visibilità fine-grained sul comportamento del sistema senza il costo di agenti user-space tradizionali. La presenza di mappe condivise, helper stabili e compilazione JIT permette di costruire pipeline di telemetria efficienti, mentre l'aggancio a hook diversi consente di intercettare eventi a livelli di astrazione differenti, dalle syscall ai punti di tracing più specifici \cite{ebpf-docs,linux-bpf-docs}.

In ambito cloud-native, \ac{ebpf} viene spesso usato come strato di osservabilità e controllo perché riduce il numero di copie dei dati e concentra la logica di raccolta vicino alla sorgente dell'evento. Questo non elimina la complessità operativa, ma sposta il problema dalla lettura passiva dei log alla progettazione di regole e filtri che siano abbastanza espressivi da distinguere attività legittime da comportamenti sospetti.

% da revisionare + aggiungere riferimenti e immagine/i

%----------------------------------------------------------------------------------------

\chapter{Falco e Tetragon}
\label{chap:falco-tetragon}
% Descrizione di Falco e Tetragon: generale e architetturale, con diagrammi di flusso dei percorsi di evento e policy.

\section{Falco}
\label{sec:falco}
Falco è una piattaforma di runtime security orientata al rilevamento. Il suo obiettivo principale è osservare il comportamento del sistema, correlare gli eventi con una regole engine e generare alert quando viene riconosciuto un pattern anomalo o potenzialmente malevolo \cite{falco-docs}.

Nell'architettura corrente Falco monitora soprattutto le syscall e, più in generale, sorgenti di eventi che possono essere estese tramite plugin. Questo approccio gli consente di raccogliere contesto sia dal kernel sia da metadati provenienti da container runtime e Kubernetes, mantenendo una separazione netta tra acquisizione dell'evento e valutazione della regola. Le regole sono espresse in YAML e vengono valutate dal demone user-space, che si occupa anche della formattazione e dell'inoltro degli alert verso stdout, syslog, webhook o altri sistemi di downstream \cite{falco-docs}.

Il vantaggio di questo modello è la flessibilità: una volta che l'evento è stato raccolto, la logica di rilevamento può essere modificata senza toccare il kernel. Di contro, il controllo è reattivo, cioè interviene dopo che l'evento è già avvenuto; per questo Falco è particolarmente utile per detection, audit e incident response, ma non per bloccare direttamente l'esecuzione di un'azione nel punto in cui avviene \cite{falco-docs}.

Dal punto di vista operativo, Falco tende a essere adottato quando serve una visibilità ampia e un'integrazione semplice con flussi di alerting esistenti. La sua efficacia dipende però dalla qualità del tuning delle regole e dalla capacità di ridurre il rumore prodotto da workload eterogenei, soprattutto in ambienti con elevata dinamicità come Kubernetes.

\section{Tetragon}
\label{sec:tetragon}
Tetragon è un progetto di Cilium pensato per security observability e runtime enforcement a livello di kernel. A differenza di un approccio puramente reattivo, Tetragon cerca di combinare raccolta eventi, filtraggio e applicazione della policy all'interno del percorso di esecuzione dell'evento, così da ridurre la finestra temporale tra osservazione e mitigazione \cite{tetragon-home}.

Il modello concettuale ruota attorno alle tracing policy: il sistema osserva processi, syscall, file e traffico di rete, correlandoli con l'identità del workload e con il contesto Kubernetes. Quando una policy viene violata, Tetragon può notificare l'evento ma anche applicare un'azione di enforcement immediata, per esempio interrompendo l'esecuzione del processo o bloccando un'operazione prima che produca effetti rilevanti \cite{tetragon-home}.

Questo approccio sfrutta il fatto che l'elaborazione avviene direttamente nel kernel tramite \ac{ebpf}, riducendo l'overhead introdotto dal passaggio continuo in user space. In pratica Tetragon è più vicino a un motore di prevenzione che a un semplice sistema di monitoraggio: il suo valore aggiunto emerge soprattutto quando la policy deve essere applicata in modo coerente e con latenza minima, ad esempio per impedire l'esecuzione di binari non consentiti o rilevare tentativi di esfiltrazione di dati \cite{tetragon-home}.

La disponibilità di metadati del kernel, incluso il supporto a informazioni di tipo come BTF, rende inoltre più robusto l'aggancio alle strutture interne e facilita l'adattamento a kernel diversi, purché l'ambiente soddisfi i requisiti richiesti dal sottosistema BPF \cite{linux-btf-docs,linux-bpf-docs}.

% TODO: DA COMPLETARE E REVISIONARE

%----------------------------------------------------------------------------------------

\chapter{Contesti d'uso e integrazione}
\label{chap:use-cases}
% Requisiti di Sistema e dipendenze del Kernel: BTF (CONFIG_DEBUG_INFO_BTF), versioni del kernel Linux e permessi CAP_SYS_ADMIN/CAP_BPF.
% Analisi e Prova nei 4 Contesti di Deployment: Standalone Container, Docker Swarm, Kubernetes Cluster, Ambienti Cloud Generici.
% Valutazione sulla Facilità di Configurazione: Complessità nella scrittura e nel tuning delle regole (Sintassi YAML di Falco vs. CRD TracingPolicy di Tetragon).

%----------------------------------------------------------------------------------------

\chapter{Progettazione Suite di Test}
\label{chap:testing}
% Infrastruttura di Laboratorio: Configurazione dell'ambiente di test
% Architettura dell'Applicazione Cavia
% Progettazione della Suite di Testing Automatizzata

%----------------------------------------------------------------------------------------

\chapter{Risultati e Discussione} % sezioni da definire meglio
\label{chap:results}
% Rilevazione ed Evasione (Detection & Bypass)
% Studio dei Falsi Positivi
% Reattività ed Enforcement (Prevention Test)
% Meccanismi di Avviso e Automazione della Risposta

%----------------------------------------------------------------------------------------

\chapter{Osservabilità: Integrazione con Grafana, Prometheus}
\label{chap:observability}
% SOLO SE AVANZA TEMPO
% Integrazione con Grafana, Prometheus

%----------------------------------------------------------------------------------------

\chapter{Conclusione e sviluppi futuri}
\label{chap:conclusion}
% Sintesi dei Risultati: Tabella comparativa finale su Efficacia, Robustezza, Usabilità e Prestazioni.
% Linee Guida di Scelta: Guida per un Ingegneer/Security Architect su quando adottare Falco, Tetragon o un approccio ibrido.
% Sviluppi Futuri


%----------------------------------------------------------------------------------------
% BIBLIOGRAPHY
%----------------------------------------------------------------------------------------

\backmatter

% \nocite{*} % Remove this as soon as you have the first citation

\bibliographystyle{plain}
\bibliography{bibliography}

\begin{acknowledgements}
Opzionale: ringraziamenti a supervisori, colleghi, amici e famiglia.
\end{acknowledgements}

\end{document}
